\documentclass[12pt]{article}

% ИСПРАВЛЕННЫЙ PREAMBLE с необходимыми пакетами
\usepackage[T2A]{fontenc}
\usepackage[utf8]{inputenc}
\usepackage[russian,english]{babel}
\usepackage{amsmath,amssymb,amsfonts}
\usepackage{mathtools}
\usepackage{bm}

% Геометрия страницы
\setlength{\textwidth}{168mm}
\setlength{\textheight}{210mm}
\topmargin -10mm
\oddsidemargin  -1mm
\evensidemargin -1mm

\begin{document}
	\title{Космологическая постоянная}
	
	\author{Бородин Никита}
	
	
	
	\maketitle
	\newpage
	\section{Введение}
	
	\section*{\textbf{Постановка задачи}}
	\begin{itemize}
		\item 	Покажите, что ненулевая (бесконечная!) вакумная плотность энергии квантового поля переопределяет космологическую
		постоянную в уравнениях Эйнштейна.
		\item  Обсудите как меняется решение эволюции однородной
		и изотропной Вселенной в зависимости от вклада в вакуумную
		энергию фермионных и бозонных полей.
		\item  Исследуйте литературу с современными экспериментальными данными по измерению космологической постоянной в
		модели однородной и изотропной Вселенной. Укажите, какими
		свойствами должна обладать система, чтобы объяснить наблюдаемое ускоренное расширение Вселенной.
		\item Рассмотрите простейшую суперсимметричную модель Весса-
		Зумино[4], содержащую один комплексный скаляр и один майорановский фермион. Получите значение средней плотности
		энергии вакуума. Решает ли суперсимметрия проблему космо-
		логической постоянной?
	\end{itemize}
	
	Проблема космологической постоянной заключается в огромном расхождении между теоретическми предсказаниями квантовой теории поля и наблюдаемыми значениями.
	
	рассмотрим стандартную форму уравнения Эйнштейна
	\[
	 R_{\mu \nu} - \frac{1}{2} g_{\mu \nu} R + \Lambda g_{\mu \nu} = \frac{8 \pi G}{c^4} T_{\mu \nu}
	\]
	
	В квантовой теории поля имеет ненулевую плотность энергии $\rho_{vac}$. Вакуум характеризуется тензором энергии-импульса:
	\[
	T_{\mu \nu}^{(vac)} = \rho_{vac} g_{\mu \nu}
	\]
	
	Общий тензор энергии-импульса включает вклады материи и вакуума:
	
	\[
	T_{\mu \nu} = T_{\mu \nu}^{(matter)} + T_{\mu \nu}^{(matter)} = T_{\mu \nu}^{(matter)} + \rho_{vac} g_{\mu \nu}
	\]
	
	После подстановки получаем
	\[
	R_{\mu \nu} - \frac{1}{2} g_{\mu \nu} + (\Lambda + \frac{8 \pi G}{c^4} \rho_{vac}) g_{\mu \nu} = \frac{8 \pi G}{c^4} T_{\mu \nu}^{(matter)} 
	\]
	
	Полученная эффективная космологическая постоянная
	\[
	\Lambda_{eff} = \Lambda + \kappa \rho_{vac}
	\]
	$\kappa \sim 2.07 \times 10^{-43}$ 
	Рассмотрим различные квантовые поправки от вакуумной энергии для разных масштаб энергии
	 
	 В КЭД
	 \[
	 E = <0| \hat{H}|0> = \frac{1}{2} <0| \int d^3 x (\hat{\mathbf{E}}^2 + \hat{\mathbf{B}}^2)|0> = \delta^3(0) \int d^3 k \frac{1}{2} \hbar \omega_{\mathbf{k}}
	 \]
	 
	 Тогда 
	 
	 \[
	 \rho_{vac} = \frac{E}{V} = \frac{1}{V} \sum_{\mathbf{k}} \frac{1}{2} \hbar \omega_{\mathbf{k}} \sim \frac{\hbar}{2 \pi^2 c^3} \int_{0}^{\omega_{max}} \omega^3 d\omega = \frac{\hbar}{8 \pi^2 c^3} \omega^4_{max}
	 \]
	 
	 \end{document}